\tzpanchor{1220}
\tzpentryheader{1220}{Sumerian Lexicon of Hand Gestures and Actions}

\begin{tzpsnapshot}
\tzpsnaprow{TZPID}{\tzpcode{ID1220}}
\tzpsnaprow{UUID}{\tzpcode{4869b0a8-1ef8-5920-965a-ff69ed388984}}
\tzpsnaprow{Type}{Dynamical / Variational Statement}
\tzpsnaprow{Scale}{field-dynamical}
\tzpsnaprow{LOC Classification}{Working routing: QC20.7 field theory; QA321 variational structure}
\tzpsnaprow{arXiv Classification}{Primary: hep-th; Cross-list: math-ph, gr-qc}
\end{tzpsnapshot}

\tzpsection{Canonical Statement}
2. The Curling / Grasping Hand (Screenshots 343, 344, 346, 348)

\[
r  # Convert spherical back to unit vector for angle measurement\r  # Note: Theta/Phi in JSON are in degrees\r  theta = np
\]
\[
radians(entry[\""spherical\""][\""theta (polar)\""])\r  phi = np
\]

\noindent
\begin{minipage}[t]{0.48\linewidth}
\tzpsection{Wolfram Form}
\begingroup
\small
\[
\begin{aligned}
\mathfrak{W}_{\mathrm{TZP}}\!\left(\mathcal{K}_{1220}\right)
&\longmapsto
\left\langle \tau(\mathcal{L}),\; \Omega^{\mathbb{T}},\; \Psi_{\mathrm{T}} \right\rangle
\end{aligned}
\]
\[
\mathcal{K}_{1220}
:=
\operatorname{HoldForm}\!\left(\mathcal{T}_{1220}\right)
\]
\[
\begin{aligned}
\operatorname{Admissible}_{\mathrm{TZP}}\!\left(\mathcal{K}_{1220}\right)
&\Longleftrightarrow
\mathfrak{W}_{\mathrm{TZP}}\!\left(\mathcal{K}_{1220}\right)\not\equiv\varnothing
\end{aligned}
\]
\textbf{Wolfram register:} canonical symbol lift\\
\textbf{Title lift:} Sumerian Lexicon of Hand Gestures and Actions\\
\textbf{Symbol key:} \(\tau(\mathcal{L})\): topological charge density\\ \(\Omega^{\mathbb{T}}\): Trawinistic boundary curvature\\ \(\Psi_{\mathrm{T}}\): Trawinistic wavefunction
\par
\endgroup
\end{minipage}\hfill
\begin{minipage}[t]{0.48\linewidth}
\tzpsection{TRAWIN Normalization}
\[
\tzpnormalizewrap{r  # Convert spherical back to unit vector for angle measurement\r  # Note: Theta/Phi in JSON are in degrees\r  theta = np,\,radians(entry[\""spherical\""][\""theta (polar)\""])\r  phi = np,\,n # y = r sin(theta) sin(phi)\r  # z = r cos(theta)\r  \r  x = np}
\]

\small
\textbf{T}: Unified Density is localized within the engineering and implementation arcs arc of the directory.\\
\textbf{R}: Euler-Lagrange Flow preserves the operative relation that sumerian lexicon of hand gestures and actions requires.\\
\textbf{A}: Stationary Closure fixes the admissible closure condition for the entry.\\
\textbf{W}: Sumerian Lexicon of Hand Gestures and Actions is rendered as a reusable operator-level statement.\\
\textbf{I}: The informational content of sumerian lexicon of hand gestures and actions is retained rather than flattened.\\
\textbf{N}: Sumerian Lexicon of Hand Gestures and Actions closes as a distinct normalized TRAWIN object.
\normalsize
\end{minipage}

\tzpsection{Derivable Consequences}
The entry supplies the book with a true unification density, so later extraction and dynamo pages can be read as sectoral consequences of one action.
\clearpage

\tzpentryheader{1220}{Oxford Dictionary of the Queens, NY English}

\begingroup\small
\tzpsection{Definition}
Sumerian Lexicon of Hand Gestures and Actions is a registry definition in Engineering and Implementation Arcs with secondary pressure from Registry and Publication Workflow.

% BEGIN TZP CONTEXTUAL MEANING
\tzpsection{Contextual Meaning}
\begingroup\footnotesize\setlength{\emergencystretch}{3em}\sloppy
\textbf{Formal statement:} 2. The Curling / Grasping Hand (Screenshots 343, 344, 346, 348). \textbf{Contextual meaning:} The entry reads Sumerian Lexicon of Hand Gestures and Actions as a controlled technical headword whose equation layer, proof route, and registry role are interpreted together. \textbf{Derivable consequences:} The entry supplies the book with a true unification density, so later extraction and dynamo pages can be read as sectoral consequences of one action. \textbf{Lemma support:} Unified Density; Euler-Lagrange Flow; Stationary Closure.\par
\endgroup
% END TZP CONTEXTUAL MEANING

\tzpsection{Technical Interpretation}
Technically, Sumerian Lexicon of Hand Gestures and Actions is read through the local equation chain and the TRAWIN normalization recorded on the entry page.

\tzpsection{Equation Grounding}
Equation anchors: r \# Convert spherical back to unit vector for angle measurementr \# Note: Theta/Phi in JSON are in degreesr theta = np; radians(entry[ spherical ][ theta (polar) ])r phi = np; and n \# y = r sin(theta) sin(phi)r \# z = r cos(theta)r r x = np. Proof route: show that the joined density yields a coherent Euler-Lagrange system whose sectoral terms remain admissible under TRAWIN normalization. Formalization route: prove that stationarity of the unified action implies the stated field equation.

\tzpsection{Interpretive Role}
The entry supplies the book with a true unification density, so later extraction and dynamo pages can be read as sectoral consequences of one action.
\endgroup

\tzpsection{TZP Thesaurus}
\begin{enumerate}[leftmargin=1.6em,itemsep=6pt,topsep=4pt]
\item \textbf{Sumerian Lexicon Hand Gestures Actions}: entry-specific alternate wording for indexing, related literature, and local formal context.
\item \textbf{Sumerian Lexicon Hand Variational Analogue}: variational synonym for action, energy generation, admissible variation, or closure.
\item \textbf{Actions Equivalence}: entry-specific alternate wording for indexing, related literature, and local formal context.
\end{enumerate}

\tzpsection{Encyclopedia Mathematica}
\begin{samepage}
\noindent\begin{minipage}[t]{0.45\linewidth}
\vspace{0pt}
\raggedright
\scriptsize\textbf{Image reading.} The rendered manifold at right is the per-ID light plate for Sumerian Lexicon of Hand Gestures and Actions. It should be read as the visual companion to the entry's equation layer, not as a repeated template diagram.\par
\end{minipage}\hfill
\begin{minipage}[t]{0.53\linewidth}
\vspace*{-0.10in}
\raggedright
\tzpencyclopediaimage{1220}
\end{minipage}
\end{samepage}

\clearpage

\tzpentryheader{1220}{Direct Equation Progression and Proof Trajectory}

\tzpsection{Direct Equation Progression}
\begin{enumerate}[leftmargin=1.6em,itemsep=9pt,topsep=6pt]
\item \textbf{Unified Density}
\[
r  # Convert spherical back to unit vector for angle measurement\r  # Note: Theta/Phi in JSON are in degrees\r  theta = np
\]
This gathers the sectors that the quaternary arc is trying to hold in one variational frame.

\item \textbf{Euler-Lagrange Flow}
\[
radians(entry[\""spherical\""][\""theta (polar)\""])\r  phi = np
\]
This turns the density into an actual field law rather than a formal sum.

\item \textbf{Stationary Closure}
\[
n # y = r sin(theta) sin(phi)\r  # z = r cos(theta)\r  \r  x = np
\]
This closes the entry by imposing stationary admissibility on the total action.
\end{enumerate}

\tzpsection{Proof}
\textbf{Proof route.} show that the joined density yields a coherent Euler-Lagrange system whose sectoral terms remain admissible under TRAWIN normalization.

\textbf{Formal method.} prove that stationarity of the unified action implies the stated field equation.

\medskip
\tzpproofartifacts{Isabelle audit: corpus classification recorded in appendix.}{Lean audit: compiled corpus certificate.}{Rocq audit: compiled corpus certificate.}

\tzpsection{Dependencies and Test Handle}
\begin{tzpsnapshot}
\tzpsnaprow{Dependencies}{\tzpidref{1219}, \tzpidref{1200}}
\tzpsnaprow{Node}{\tzpcode{registry_id_1220}}
\tzpsnaprow{Support Titles}{Unified Density; Euler-Lagrange Flow; Stationary Closure}
\tzpsnaprow{Test Handle}{If the combined density does not produce a coherent stationary variational law, the unification entry fails.}
\end{tzpsnapshot}

\tzpsection{Canonical Equation Links}
\tzpequationlinks
  {\mathcal{L}_{U} = \cdots}
  {\frac{\partial \mathcal{L}_{U}}{\partial \Phi} - \partial_{\mu}\!\left(\frac{\partial \mathcal{L}_{U}}{\partial(\partial_{\mu}\Phi)}\right) = \cdots}
  {n # y = r sin(theta) sin(phi)\r  # z = r cos(theta)\r  \r  x = np}
